\documentclass[11pt]{article}
\usepackage[margin=25mm]{geometry}
\usepackage{fontspec}
\setmainfont{Latin Modern Roman}
\usepackage{amsmath,amssymb,mathtools}
\usepackage{unicode-math}
\setmathfont{Latin Modern Math}
\usepackage{xurl,hyperref}
\hypersetup{colorlinks=true,urlcolor=blue}
\setcounter{secnumdepth}{0}
\setlength{\parindent}{0pt}
\setlength{\parskip}{5pt}
\setlength{\emergencystretch}{3em}
\newcommand{\sref}[2]{\hyperlink{source:#1:#2}{[S#2]}}
\newcommand{\cataloguescope}[1]{\paragraph{Recorded target.}#1}
\newcommand{\R}{\mathbb R}
\newcommand{\C}{\mathbb C}
\newcommand{\N}{\mathbb N}
\newcommand{\Z}{\mathbb Z}
\newcommand{\Q}{\mathbb Q}
\newcommand{\F}{\mathbb F}
\newcommand{\eps}{\varepsilon}
\begin{document}
\section{A two-vertex counterexample to color-number monotonicity}\label{statement}

\subsection{Definitions and mathematical statement}

\paragraph{Shared notation.}$\mathbb N=\{1,2,\ldots\}$; the standard number fields have their usual meanings. All additional conventions are stated locally.


\paragraph{Definitions and notation.}For a finite graph $G$, a proper $q$-coloring assigns colors to vertices with adjacent colors distinct. Single-site Glauber dynamics selects a vertex and color uniformly, accepting the recoloring only when proper. Its mixing time is $t(q)=\min\{t:\max_\sigma\|P_q^t(\sigma,\cdot)-\pi_q\|_{\mathrm{TV}}\le1/e\}$, where $\pi_q$ is uniform on proper colorings. Random $k$-SAT has $n$ variables and $m=rn$ clauses; a clustering transition concerns the geometry of its solution set.

\paragraph{Catalogue wording retained for traceability.}
The record includes source questions (5)--(16), not only the coloring question.
\begin{enumerate}
\item For a finite graph of maximum degree $\Delta$, does $t(q)$ decrease as $q$ increases, in the ergodic range $q\ge\Delta+2$? For a square lattice with $n$ sites, determine mixing behavior for $q=4,5$, in the source’s boundary convention; it records historical polynomial behavior for $q=3$ and $O(n\log n)$ for $q\ge6$.
\item Is some $k$-SAT variant, possibly on a lattice, a suitable Boolean model for the scaling of jamming?
\item Construct a constrained phase-transition model, soften its constraints and determine sensitivity to that change.
\item Identify useful response or correlation functions for constraint-satisfaction problems such as $k$-SAT and determine their transition singularities or power laws.
\item Explain the physical meaning, if any, of the finite external field used in the refined second-moment satisfiability bound, including whether it minimizes suitable zero-temperature free-energy fluctuations.
\item Find algorithms that empirically solve in roughly linear time at values of $k$ and clause density where clustering has been rigorously established.
\item Prove clustering below the satisfiability transition for $k=3,4,5$, in the source’s setting.
\item Prove an order/disorder transition for equilibrium hard spheres or hard disks.
\item Determine how to incorporate friction into a jamming model.
\item For Ising Glauber dynamics on the $d$-dimensional lattice at fixed $\beta$ (and possibly fixed external field), determine whether sharp total-variation cutoff in the large-volume limit occurs exactly when the infinite-volume Gibbs state is unique.
\item Define a mathematically useful thermodynamic limit of random $k$-SAT with $n\to\infty$ and $m=rn$ at constant $r$.
\item For an Ising model in a finite $\mathbb Z^d$ box with $\beta>\beta_c$, zero field and plus boundary, determine whether Glauber mixing from the source’s warm starts (all plus or random) is $O(n\log n)$, with $n$ sites.
\end{enumerate}
\paragraph{Scope and source qualification.}The inherited solved label and title refer to a coloring subcase in the research metadata and are not a verified solution of this twelve-question bundle. All targets and the source’s historical partial context are retained; no current-status claim is inferred from the old workshop list. The catalogue title/status was paired with a machine-generated research attempt; that attempt is not peer-reviewed evidence of settlement. The statement covers the original source questions rather than only the title’s subcase.
\subsection{Short English statement}
Study coloring dynamics, constraint-satisfaction transitions, jamming, hard-particle ordering and Ising mixing.
\subsection{Research attempt}
\paragraph{Model, date and investigated scope.}
OpenAI GPT-6 Astra (configured model \texttt{gpt-6-astra}), ultra reasoning; generated 11 October 2026. Reconstructed a two-isolated-vertex counterexample t(2)=1<t(3)=2 to color-number mixing monotonicity. The other eleven clauses and square-lattice cases remain unresolved; source-solved metadata is preserved.

\paragraph{Sources and literature check.}
Read the supplied complete definition, catalogue wording, scope qualifications and source list for this canonical ID. The original-source citations are inherited as supplied references; no fresh primary-source or novelty search was performed for this attempt. No current open-status claim is inferred from this calculation.

\paragraph{An exact two-vertex coloring-chain calculation.}Take a graph consisting of two isolated vertices, so \(\Delta=0\) and every \(q\ge2\) is in the stated ergodic range. At each discrete step, one vertex is chosen uniformly and its color is replaced by a uniform color; every proposal is accepted. Start at \((1,1)\). After \(t\ge1\), the probability that one particular vertex has never been selected is \(2^{-t}\), and both cannot be unselected. Conditional on both having been selected, the final colors are independent uniform; conditional on only one being selected, the other remains one. Therefore
\[P^t-\pi=2^{-t}(\delta_1\otimes u+u\otimes\delta_1-2u\otimes u),
\]
where \(u\) is the uniform measure on \(q\) colors. Summing absolute values over the cases where zero, one or two coordinates equal one gives
\[\|P^t-\pi\|_{\mathrm{TV}}=2^{-t}\,2(q-1)^2/q^2.
\]
All starting states have the same distance by color symmetry. At \(q=2\), the distance at one step is \(1/4<1/e\), while at time zero it is \(3/4\), so \(t(2)=1\). At \(q=3\), the one-step distance is \(4/9>1/e\) and the two-step distance is \(2/9<1/e\), so \(t(3)=2\). Thus monotonic decrease fails for this explicit finite graph under the source's threshold and update rule.
\paragraph{Scope of the reconstruction.}This elementary counterexample concerns only the universal color-number monotonicity clause. It does not answer square-lattice four/five-color mixing, Boolean jamming models, soft constraints, CSP responses or field interpretation, algorithmic clustering questions, small-k clustering, hard-particle ordering, friction, cutoff-versus-uniqueness, random-SAT thermodynamic limits or warm-start Ising mixing. The inherited source-solved label refers to a coloring subcase and is not promoted to a solution of the twelve-question record. This reconstruction is not claimed novel, and the complete record remains unresolved in this attempt.

\paragraph{Verification and adversarial audit.}
Checked maximum degree zero, ergodic q range, exact mixture law, all TV sign classes, integer threshold times and full twelve-question scope. This was an LLM adversarial audit, not independent human verification; no theorem prover was used. Prompts A then B applied to this ID: Checked maximum degree zero, ergodic q range, exact mixture law, all TV sign classes, integer threshold times and full twelve-question scope. Scope remains limited: The concrete counterexample settles only the monotonicity clause; all other bundled models and transition/mixing targets need independent arguments.

Coordinator checked the exact two-vertex transition mixture, total-variation sum and threshold mixing times at two and three colors; only the specified clause is contradicted. This additional check was model-assisted, not independent human verification.

\paragraph{Tools, process and contributions.}
Python 3 and local shell used to read supplied files and serialize this manuscript; mathematics derived symbolically in the text, with no numerical experiment or formal proof assistant. Contribution prompt A was followed by prompt B. GPT-6 Astra with ultra reasoning drafted the research and performed its stated self-audit; any additional coordinator contribution is identified in the verification and audit. A GPT-6-based Codex coordinator assembled the manuscript and checked repository integration; its exact variant and reasoning configuration were not exposed. No human mathematical contribution, independent human proof check or formal proof certification is supplied. Complete derivations and the consulted source references are included here; private reasoning transcripts are not published.

\paragraph{Final status and first remaining gap.}
\textbf{UNRESOLVED}. The concrete counterexample settles only the monotonicity clause; all other bundled models and transition/mixing targets need independent arguments.

\section{Completion Estimate}
\noindent\textbf{COMPLETION: 0\%}
The calculation is an elementary benchmark or conditional reduction; it establishes no claimed new portion of the complete source target.

\subsection{Sources}
\begin{itemize}
\item[\hypertarget{source:statement:1}{S1}] ULAM.AI and the UnsolvedMath Contributors, UnsolvedMath, catalogue entry. Catalogue curation under CC BY 4.0; definitions and English exposition added locally.
\url{https://huggingface.co/datasets/ulamai/UnsolvedMath}.
\item[\hypertarget{source:statement:2}{S2}] American Institute of Mathematics, Phase transitions, Questions 5–16.
\url{https://aimath.org/WWN/phasetransition/phasetransition.pdf}.
\end{itemize}
\label{end:statement}
\end{document}
